\exercisesection{\S~1}

\begin{enumerate}
\def\labelenumi{\arabic{enumi})}
\tightlist
\item
  Let \(\mathrm{G}\) be a finite dimensional, connected, commutative complex Lie group, and let \(\mathrm{V}\) be its Lie algebra.
\end{enumerate}

\begin{enumerate}
\def\labelenumi{\alph{enumi})}
\item
  The map \(\exp_{\mathrm{G}}: \mathrm{V} \to \mathrm{G}\) is a surjective homomorphism; its kernel \(\Gamma\) is a discrete subgroup of \(\mathrm{V}\) .
\item
  G is compact if and only if \(\Gamma\) is a lattice in V; then G is said to be a complex torus.
\item
  Let \(\Gamma\) be the discrete subgroup of \(\mathbf{C}^2\) generated by the elements \(e_1 = (1,0), e_2 = (0,1), e_3 = (\sqrt{2},i)\) ; put \(\mathrm{G} = \mathrm{C}^2 / \Gamma\) , \(\mathrm{H} = (\Gamma + \mathrm{C}e_1) / \Gamma\) . Show that H is isomorphic to \(\mathbf{C}^*\) , that G/H is a complex torus of dimension 1, but that G contains no non-zero complex torus.
\item
  Every finite dimensional connected compact complex Lie group is a complex torus (cf.~Chap. III, §6, no. 3, Prop. 6).
\end{enumerate}

\begin{enumerate}
\def\labelenumi{\arabic{enumi})}
\setcounter{enumi}{1}
\tightlist
\item
  Let H be the set of complex numbers \(\tau\) such that \(\mathcal{I}(\tau) > 0\) .
\end{enumerate}

\begin{enumerate}
\def\labelenumi{\alph{enumi})}
\item
  Show that an analytic law of left operation of the discrete group \(\mathbf{SL}(2,\mathbf{Z})\) on H can be defined by putting \(\gamma\tau=\frac{a\tau+b}{c\tau+d}\) for all \(\gamma=\begin{pmatrix}a&b\\ c&d\end{pmatrix}\in\mathbf{SL}(2,\mathbf{Z})\) and \(\tau\in H\) .
\item
  For \(\tau\in H\) , denote by \(T_{\tau}\) the complex torus \(\mathbf{C}/(\mathbf{Z}+\mathbf{Z}\tau)\) . Show that the map \(\tau\mapsto T_{\tau}\) induces by passage to the quotient a bijective map from the set \(\mathrm{H}/\mathrm{SL}(2,\mathbf{Z})\) to the set of isomorphism classes of complex tori of dimension one.
\end{enumerate}

\begin{enumerate}
\def\labelenumi{\arabic{enumi})}
\setcounter{enumi}{2}
\item
  Let G be an integral subgroup of \(\mathbf{O}(n)\) such that the identity representation of G on \(\mathbf{R}^{n}\) is irreducible. Show that G is closed (write G in the form \(K \times N\) , where K is compact and N is commutative, and prove that \(\dim \overline{N} \leq 1\) ).
\item
  Show that the conditions in Prop. 3 (no. 3) are equivalent to each of the following conditions:
\end{enumerate}

(ii') The group \(\operatorname{Ad}(\mathbf{G})\) is relatively compact in \(\operatorname{Aut}(\mathbf{L}(\mathbf{G}))\) .

(ii'\,') The group \(\operatorname{Ad}(\mathbf{G})\) is relatively compact in \(\operatorname{End}(\mathbf{L}(\mathbf{G}))\) .

\begin{enumerate}
\def\labelenumi{(\alph{enumi})}
\setcounter{enumi}{21}
\tightlist
\item
  Every neighbourhood of the identity element e of G contains a neighbourhood of e stable under inner automorphisms (cf.~Integration, §3, no. 1, Prop. 1).
\end{enumerate}

\begin{enumerate}
\def\labelenumi{\arabic{enumi})}
\setcounter{enumi}{4}
\tightlist
\item
  Let A be the closed subgroup of \(\mathbf{GL}(3,\mathbf{R})\) consisting of the matrices \((a_{ij})\) such that \(a_{ij}=0\) for i\textgreater j, \(a_{ii}=1\) for \(1\leq i\leq3\) , \(a_{12}\in \mathbf{Z}\) , \(a_{23}\in \mathbf{Z}\) ; let \(\theta\in \mathbf{R}-\mathbf{Q}\) and let B be the subgroup of A consisting of the matrices \((a_{ij})\) such that \(a_{12}=a_{23}=0\) and \(a_{13}\in\theta\mathbf{Z}\) , and let G=A/B.
\end{enumerate}

\begin{enumerate}
\def\labelenumi{\alph{enumi})}
\item
  Show that G is a Lie group of dimension one, whose Lie algebra is compact.
\item
  Show that \(\mathrm{C}(\mathrm{G}) = \overline{\mathrm{D}(\mathrm{G})} = \mathrm{G}_{0}\) , and that \(G_{0}\) is the largest compact subgroup of G.
\item
  Show that G is not the semi-direct product of \(G_{0}\) with any subgroup.
\end{enumerate}

\begin{enumerate}
\def\labelenumi{\arabic{enumi})}
\setcounter{enumi}{5}
\item
  Show that a (real) Lie algebra is compact if and only if it admits a basis \((e_{\lambda})_{\lambda\in\mathrm{L}}\) for which the structure constants \(\gamma_{\lambda\mu\nu}\) are anti-symmetric in \(\lambda,\mu,\nu\) (that is, they satisfy \(\gamma_{\lambda\mu\nu}=-\gamma_{\mu\lambda\nu}=-\gamma_{\lambda\nu\mu}\) ).
\item
  An involutive Lie algebra is a (real) Lie algebra \(\mathfrak{g}\) equipped with an automorphism s such that \(s \circ s = 1_{\mathfrak{g}}\) . Denote by \(\mathfrak{g}^{+}\) (resp. \(\mathfrak{g}^{-}\) ) the eigenspace of \(\mathfrak{g}\) relative to the eigenvalue +1 (resp. -1) of s.
\end{enumerate}

\begin{enumerate}
\def\labelenumi{\alph{enumi})}
\item
  Show that \(\mathfrak{g}^{+}\) is a subalgebra of \(\mathfrak{g}\) and that \(\mathfrak{g}^{-}\) is a \(\mathfrak{g}^{+}\) -module; we have \([\mathfrak{g}^{-},\mathfrak{g}^{-}] \subset \mathfrak{g}^{+}\) , and \(\mathfrak{g}^{+}\) and \(\mathfrak{g}^{-}\) are orthogonal with respect to the Killing form.
\item
  Show that the following conditions are equivalent:
\end{enumerate}

\begin{enumerate}
\def\labelenumi{(\roman{enumi})}
\item
  The \(\mathfrak{g}^{+}\) -module \(\mathfrak{g}^{-}\) is simple.
\item
  The algebra \(\mathfrak{g}^{+}\) is maximal among the subalgebras of \(\mathfrak{g}\) distinct from \(\mathfrak{g}\). If these conditions are satisfied, and if \(\mathfrak{g}^{+}\) contains no non-zero ideal of \(\mathfrak{g}\), the involutive Lie algebra \((\mathfrak{g}, s)\) is said to be irreducible.
\end{enumerate}

\begin{enumerate}
\def\labelenumi{\alph{enumi})}
\setcounter{enumi}{2}
\tightlist
\item
  Assume that \(\mathfrak{g}\) is semi-simple. Show that the following conditions are equivalent:
\end{enumerate}

\begin{enumerate}
\def\labelenumi{(\roman{enumi})}
\item
  The only ideals of \(\mathfrak{g}\) stable under \(s\) are \(\{0\}\) and \(\mathfrak{g}\) ;
\item
  \(\mathfrak{g}\) is either simple or the sum of two simple ideals interchanged by \(s\) .
\end{enumerate}

Show that these conditions are satisfied when \((\mathfrak{g}, s)\) is irreducible (observe that \(\mathfrak{g}\) is the direct sum of ideals stable under s and satisfying (ii)).

\begin{enumerate}
\def\labelenumi{\alph{enumi})}
\setcounter{enumi}{3}
\item
  Assume from now on that \(\mathfrak{g}\) is compact semi-simple. Show that the involutive Lie algebra \((\mathfrak{g}, s)\) is irreducible if and only if s is different from the identity and \((\mathfrak{g}, s)\) satisfies the equivalent conditions in c) (let p be a \(\mathfrak{g}^{+}\) -submodule of \(\mathfrak{g}^{-}\) and q its orthogonal complement with respect to the Killing form; observe that \([p, q] = 0\) and deduce that \(p + [p, p]\) is an ideal of \(\mathfrak{g}\)).
\item
  Prove that \(\mathfrak{g}\) is the direct sum of a family \((\mathfrak{g}_{i})_{0\leq i\leq n}\) of ideals stable under s, such that \(\mathfrak{g}_{0}\) is fixed by s and such that the involutive algebras \((\mathfrak{g}_{i},s|\mathfrak{g}_{i})\) are irreducible for \(1\leq i\leq n\) .
\end{enumerate}

\begin{enumerate}
\def\labelenumi{\arabic{enumi})}
\setcounter{enumi}{7}
\tightlist
\item
  Let G be a compact semi-simple Lie group, u an automorphism of G of order two, K the identity component of the set of fixed points of u, and X the manifold G/K (the homogeneous space X is then said to be symmetric).
\end{enumerate}

\begin{enumerate}
\def\labelenumi{\alph{enumi})}
\item
  Show that, if G is almost simple, then K is maximal among the connected closed subgroups of G distinct from G; in other words (Chap. III, §3, Exerc. 8), the operation of G on X is primitive. The symmetric space X is then said to be irreducible.
\item
  Assume that the manifold X is simply-connected; show that it is then isomorphic to a product of manifolds each of which is isomorphic either to a Lie group or to an irreducible symmetric space.
\end{enumerate}

\begin{enumerate}
\def\labelenumi{\arabic{enumi})}
\setcounter{enumi}{8}
\tightlist
\item
  Let \(\mathfrak{a}\) be a real or complex Lie algebra, and let \(G\) be a compact subgroup of \(\operatorname{Aut}(\mathfrak{a})\) . Prove that \(\mathfrak{a}\) has a Levi subalgebra (Chap. I, §6, no. 8, Def. 7) stable under \(G\) (reduce to the case in which the radical of \(\mathfrak{a}\) is abelian, and use Integration, Chap. VII, §3, no. 2, Lemma 2).
\end{enumerate}

